% Part: second-order-logic
% Chapter: syntax-and-semantics
% Section: language-of-sol

\documentclass[../../../include/open-logic-section]{subfiles}

\begin{document}

\olfileid{sol}{syn}{lan}

\olsection{দ্বিতীয়-ক্রমের যুক্তিবিদ্যার ভাষা}

প্রথম-ক্রমের যুক্তিবিদ্যার মতো দ্বিতীয়-ক্রমের যুক্তিবিদ্যার রাশিগুলিও
\emph{!!{variable}s}, \emph{!!{constant}s}, \emph{!!{predicate}s}
এবং কখনও \emph{!!{function}s} নিয়ে গঠিত মৌলিক শব্দভান্ডার থেকে
তৈরি হয়। এগুলির সঙ্গে যৌক্তিক সংযোজক, পরিমাণসূচক এবং বন্ধনী ও
কমার মতো যতিচিহ্ন মিলিয়ে \emph{পদ} ও \emph{!!{formula}s} গঠন করা
হয়। পার্থক্য হল, বস্তু-চলরাশি ছাড়াও দ্বিতীয়-ক্রমের যুক্তিবিদ্যায়
সম্পর্ক ও অপেক্ষকের চলরাশি থাকে, এবং তাদের উপর পরিমাণায়ন করা যায়।
তাই প্রথম-ক্রমের যুক্তিবিদ্যার যৌক্তিক সংকেতগুলির সঙ্গে
দ্বিতীয়-ক্রমের যুক্তিবিদ্যায় নিচের সংকেতগুলিও থাকে:

\begin{enumerate}
\item প্রত্যেক স্থানসংখ্যা~$n$-এর দ্বিতীয়-ক্রমীয় সম্পর্ক-!!{variable}s-এর
  একটি !!{denumerable}s সেট:
  $\Obj V_0^n$, $\Obj V_1^n$, $\Obj V_2^n$, \dots
\item দ্বিতীয়-ক্রমীয় অপেক্ষক-!!{variable}s-এর একটি !!{denumerable}s সেট:
  $\Obj u_0^n$, $\Obj u_1^n$, $\Obj u_2^n$, \dots
\end{enumerate}

প্রথম-ক্রমের যুক্তিবিদ্যায় সাধারণত !!{identity}~$\eq$ অন্তর্ভুক্ত
থাকে। আকর্ষণীয় কিছু প্রকাশ করতে প্রথম-ক্রমীয় ভাষা~$\Lang{L}$-র
যুক্তিবহির্ভূত সংকেতগুলি অত্যন্ত প্রয়োজনীয়। অবশ্য যুক্তিবহির্ভূত
সংকেত ছাড়াও !!{sentences} আছে, কিন্তু শুধু~$\eq$ দিয়ে আকর্ষণীয়
কিছু বলা কঠিন। দ্বিতীয়-ক্রমের যুক্তিবিদ্যায় সম্পর্ক ও অপেক্ষকের
চলরাশির সীমাহীন জোগান থাকে বলে বিশেষ কোনো যুক্তিবহির্ভূত সংকেতের
জোগান ছাড়াই প্রথম-ক্রমীয় ভাষায় প্রকাশযোগ্য যেকোনো কথা বলা যায়।

\end{document}
